\begingroup
\setlength\LTleft{0\linewidth}
\setlength\LTright{0\linewidth}
\fontsize{12.0pt}{14.4pt}\selectfont

\setlength{\LTpost}{0mm}
\begin{longtable}{@{\extracolsep{\fill}}lrrrrrr}

\toprule
 & \multicolumn{3}{c}{\textbf{Ponderação igual}} & \multicolumn{3}{c}{\textbf{Ponderado por preço}} \\ 
\cmidrule(lr){2-4} \cmidrule(lr){5-7}
\textbf{Dia} & \textbf{Média} & \textbf{t} & \textbf{p} (1) & \textbf{Média} & \textbf{t} & \textbf{p} (1) \\ 
\midrule\addlinespace[2.5pt]
Segunda & 1,03\% & 14,87 & 0,000 & 0,25\% & 3,87 & 0,000 \\
Terça & 0,93\% & 12,95 & 0,000 & 0,20\% & 3,57 & 0,000 \\
Quarta & 0,96\% & 11,12 & 0,000 & 0,27\% & 4,10 & 0,000 \\
Quinta & 0,89\% & 8,02 & 0,000 & 0,23\% & 3,14 & 0,002 \\
Sexta & 0,98\% & 10,44 & 0,000 & 0,22\% & 2,76 & 0,006 \\
Sábado & 1,01\% & 9,39 & 0,000 & 0,37\% & 4,01 & 0,000 \\
Domingo & 1,21\% & 16,17 & 0,000 & 0,37\% & 6,05 & 0,000 \\
\bottomrule
\end{longtable}
\begin{minipage}{\linewidth}
Fonte: Elaboração própria.\\
Média do retorno de mercado em cada dia da semana, com erros-padrão de Newey-West. As defasagens seguem a regra apresentada por Wooldridge. Teste conjunto de médias iguais: p = 0,014 (ponderação igual) e p = 0,139 (ponderado por preço).\\
(1) O p-valor testa se a média do dia é zero.\\
\end{minipage}
\endgroup
